% Source: https://artofproblemsolving.com/community/c3h2050265_division_problem
% Stable identity: aops-topic-2050265-q01
\textbf{$\dfrac{5}{0}$.} We would need a number $c$ with $0\cdot c=5$. But $0\cdot c=0$ for every real $c$, so no such $c$ exists. The quotient is undefined because there is \emph{no} candidate.

\textbf{$\dfrac{0}{0}$.} We would need $c$ with $0\cdot c=0$. Every real number $c$ satisfies this, so there is no \emph{unique} such number. For example, $c=1$ and $c=7$ both work, and if $\dfrac{0}{0}$ were assigned any value we could "prove" $1=7$. The quotient is undefined because there are \emph{too many} candidates.

So $\dfrac{5}{0}$ fails existence, while $\boxed{\dfrac{0}{0}\text{ fails uniqueness}}$; in both cases no value can be assigned consistently with the rules of arithmetic.
